\section{Návrh riešenia}
V tejto sekcii sa budeme venovať samotnému návrhu riešenia našej
aplikácie začínajúc od funkčných a nefunkčných požiadaviek, cez grafické návrhy a diagrami až po samotnú implemntáciu a testovanie

\subsubsection{Funkcionálne a Nefunkcionálne požiadavky}
Definivanie funkcionálnych a nefunkcionálnych požiadaviek je podstatným krokom návrhu riešenia,
pretože nám pomáha jasne stanoviť správanie aplikácie ktoré by splnili jej cieľ

\subsubsection{Funkcionálne požiadavky}

Funkcionálne požiadavky definujú správanie systému a služby, ktoré používateľ od systému očakáva. Ich hlavným cieľom je:
\begin{enumerate}
    \item Definovať funkcie, ktoré sa od systému očakávajú.
    \item Napomôcť k lepšiemu pochopeniu funkcionality systému.
    \item Identifikovať, čo systém musí alebo nesmie vykonávať.
    \item Umožniť fungovanie systému (funkčnosť nie je priamo závislá od nefunkcionálnych požiadaviek).
    \item Zabezpečiť, že systém bude spĺňať reálne potreby používateľov.
    \item Definovať prvky nevyhnutné pre samotný chod systému.
    \item Stanoviť povinné body, ktoré musia byť v systéme implementované.
    \item Zdokumentovať požiadavky, ktoré sú zvyčajne definované priamo používateľom.
    \item Umožniť vizualizáciu a pochopenie procesov pomocou prípadov použitia (\textit{Use Case}).
\end{enumerate}

\noindent Medzi konkrétne funkcionálne požiadavky navrhovanej aplikácie patrí:
\begin{itemize}
    \item registrácia a prihlásenie (login) používateľov,
    \item verifikácia používateľa pri prihlasovaní,
    \item možnosť vytvorenia novej udalosti,
    \item pridávanie podrobností k udalosti (fotografie, popis, miesto konania, čas, typ udalosti, maximálny počet účastníkov a pod.),
    \item prehliadanie dostupných udalostí na interaktívnej mape,
    \item nastavenie polomeru (kruhu) pre zobrazovanie udalostí v okolí používateľa,
    \item filtrovanie udalostí podľa preferencií,
    \item správa účastníkov (možnosť odstrániť účastníka z vlastnej udalosti),
    \item ukladanie vybraných udalostí do zoznamu obľúbených,
    \item zobrazenie histórie (navštívené, odporúčané a nastávajúce udalosti),
    \item možnosť potvrdenia účasti na udalosti,
    \item komunikácia s ostatnými účastníkmi podujatia,
    \item vzájomné hodnotenie udalostí a ich účastníkov,
    \item zdieľanie (\textit{sharing}) podujatí,
    \item prehliadanie profilov ostatných používateľov,
    \item správa a nastavenie notifikácií,
    \item systém nahlasovania nevhodných používateľov alebo udalostí,
    \item možnosť zmeny jazyka v prostredí aplikácie,
    \item automatizovaný poradovník náhradníkov (\textit{Waitlist}),
    \item verifikácia príchodu účastníka na miesto konania udalosti,
    \item systém penalizácie alebo zablokovania používateľa, ktorý potvrdil účasť, nedostavil sa a včas sa neodhlásil.
\end{itemize}


\subsubsection{Nefunkcionálne požiadavky}

Nefunkcionálne požiadavky definujú kvalitatívne vlastnosti systému a obmedzenia, za ktorých musí systém pracovať. Ich charakteristika je nasledovná:
\begin{enumerate}
    \item Vysvetľujú vlastnosti produktu a jeho atribúty kvality.
    \item Napomáhajú k pochopeniu výkonu a správania systému.
    \item Identifikujú spôsob, akým má systém stanovené úlohy vykonávať.
    \item Samotné nefunkcionálne požiadavky bez funkčných nepostačujú na chod systému.
    \item Zabezpečujú, že systém bude spĺňať špecifické očakávania používateľov na kvalitu.
    \item Sú žiaduce pre celkový zážitok, hoci v niektorých prípadoch nie sú kritické pre základnú funkcionalitu.
    \item Ich splnenie zvyšuje konkurencieschopnosť a stabilitu softvéru.
    \item Zvyčajne ich definujú softvéroví inžinieri, vývojári alebo architekti na základe technických potrieb.
    \item Možno ich dokumentovať a chápať ako atribúty kvality (napr. bezpečnosť, rýchlosť).
\end{enumerate}

\noindent Pre navrhovanú aplikáciu boli definované tieto konkrétne nefunkcionálne požiadavky:
\begin{itemize}
    \item \textbf{Intuitívnosť:} Používateľské rozhranie musí byť navrhnuté tak, aby bola navigácia v aplikácii jasná aj bez dodatočnej dokumentácie.
    \item \textbf{Responzivita:} Rozhranie aplikácie sa musí korektne prispôsobovať rôznym rozlíšeniam a pomerom strán Android zariadení.
    \item \textbf{Rýchlosť odozvy:} Systém musí zabezpečiť rýchlu odozvu pri načítavaní udalostí z mapy a pri spracovaní dát v odporúčacom algoritme (do 2 sekúnd).
    \item \textbf{Stabilita a funkčnosť:} Systém musí garantovať spoľahlivé prihlásenie používateľa a stabilitu procesov na pozadí.
    \item \textbf{Spracovanie a zaznamenávanie dát:} Efektívne ukladanie dát do databázy Firebase Firestore so zreteľom na minimalizáciu dátových prenosov.
    \item \textbf{Zaznamenávanie polohy:} Presné a šetrné spracovanie GPS súradníc používateľa pre potreby mapy a \textit{Check-inu}.
    \item \textbf{Správa povolení:} Systém musí korektne vyžadovať a spracovávať povolenia od používateľa (prístup k polohe, galérii a kalendáru).
    \item \textbf{Ochrana údajov:} Zabezpečenie osobných údajov používateľov a šifrovaná komunikácia medzi klientom a serverom.
    \item \textbf{Kompatibilita:} Možnosť inštalácie a synchrónneho používania na viacerých Android zariadeniach pod jedným účtom.
    \item \textbf{Rozšíriteľnosť:} Architektúra kódu musí umožňovať jednoduché pridávanie nových funkcií v budúcnosti bez nutnosti prepisovania jadra aplikácie.
\end{itemize}